\documentclass[preview]{standalone}
\usepackage{amsmath}
\usepackage{amssymb}
\usepackage{stellar}
\usepackage{definitions}
\usepackage{bettelini}
\begin{document}
\id{diffeq-linear-first-order}
\genpage
\section{Linear first-order equations}

\begin{snippetdefinition}{diffeq-linear-first-order-definition}{Linear first-order equation}
    Let \(I\subseteq\realnumbers\) be an open interval and \(p,q\in C(I)\).
    A \emph{linear first-order equation in normal form} is
    \(y'+p(x)y=q(x)\), equivalently \(y'=q(x)-p(x)y\).
\end{snippetdefinition}

\begin{snippettheorem}{linear-first-order-differential-equation-solution}{Integrating-factor formula}
    Let \(p,q\in C(I)\) on an open interval \(I\), and fix \(x_0\in I\).
    Put \(P(x)=\int_{x_0}^x p(t)\,dt\). All solutions of \(y'+py=q\) on \(I\) are
    \[
        y(x)=e^{-P(x)}\left(C+\integral[x_0][x][e^{P(t)}q(t)][t]\right),
        \qquad C\in\realnumbers.
    \]
    The initial condition \(y(x_0)=y_0\) selects \(C=y_0\), giving a unique
    solution on all of \(I\).
    Equivalently, for \(y'=a(x)y+b(x)\) and a primitive \(A'=a\),
    \[
        y(x)=e^{A(x)}\left(C+\int e^{-A(x)}b(x)\,dx\right),
    \]
    with any one fixed primitive in the last expression.
\end{snippettheorem}

\begin{snippetproof}{linear-first-order-differential-equation-solution-proof}{linear-first-order-differential-equation-solution}{Integrating-factor formula}
    For the homogeneous equation \(y'=ay\), division by a nonzero solution gives
    \[
        \frac{y'}{y}=a,\qquad \ln|y|=A+C_0,\qquad y=Ce^A.
    \]
    The missing zero solution is included by \(C=0\).
    For \(y'+py=q\), seek a nonzero integrating factor \(\mu\) such that
    \[
        \mu y'+p\mu y=(\mu y)'=\mu y'+\mu' y.
    \]
    Thus \(\mu'=p\mu\); choosing \(\mu=e^P\) gives \(\mu(x_0)=1\).
    Multiplication and integration now yield
    \begin{align*}
        (\mu y)'&=\mu q,\\
        \mu(x)y(x)-y(x_0)&=\integral[x_0][x][\mu(t)q(t)][t],\\
        y(x)&=e^{-P(x)}
        \left(y(x_0)+\integral[x_0][x][e^{P(t)}q(t)][t]\right).
    \end{align*}
    This proves necessity and uniqueness. Conversely, the displayed expression
    is defined and continuously differentiable throughout \(I\); differentiating
    verifies the equation, and evaluation at \(x_0\) verifies the initial value.
    Taking \(p=-a\) recovers the alternative notation, with integrating factor \(e^{-A}\).
\end{snippetproof}

\section{Examples}

\begin{snippetexample}{diffeq-linear-gaussian-example}{Polynomial forcing}
    Solve \(y'+2xy=x^3\), \(y(0)=1\).
    The coefficients are continuous on \(\realnumbers\) and \(P(x)=x^2\), so
    \[
        y(x)=e^{-x^2}\left(1+\integral[0][x][t^3e^{t^2}][t]\right).
    \]
    With \(s=t^2\), integration by parts gives
    \[
        \int t^3e^{t^2}\,dt=\frac12(t^2-1)e^{t^2},
        \qquad
        \integral[0][x][t^3e^{t^2}][t]
        =\frac12(x^2-1)e^{x^2}+\frac12.
    \]
    Consequently the unique global solution is
    \[
        y(x)=e^{-x^2}\left(\frac32+\frac12(x^2-1)e^{x^2}\right)
        =\frac32e^{-x^2}+\frac{x^2-1}{2}.
    \]
\end{snippetexample}

\begin{snippetexample}{diffeq-linear-tangent-example}{Choosing the coefficient interval}
    Consider \(y'+y\tan x=3\sin x\), \(y(\pi)=1\).
    The coefficient \(\tan x\) is continuous on
    \((-\pi/2+k\pi,\pi/2+k\pi)\). The interval containing the initial point is
    \(I=(\pi/2,3\pi/2)\). There,
    \[
        P(x)=\integral[\pi][x][\tan t][t]=-\ln|\cos x|,
        \qquad e^{-P(x)}=|\cos x|=-\cos x.
    \]
    The integrating factor is \(1/|\cos x|\), so
    \begin{align*}
        y(x)&=-\cos x\left(1+
        3\integral[\pi][x][\frac{\sin t}{|\cos t|}][t]\right)\\
        &=-\cos x\left(1+3\ln(-\cos x)\right).
    \end{align*}
    The absolute value matters: \(\cos x<0\) throughout the chosen interval.
\end{snippetexample}
\end{document}
